\documentclass{pgnotes}

\title{Storage}

\begin{document}

\maketitle

\section{Key storage questions}

For our particular application / use case: 
\begin{enumerate}
\item \textbf{How much} storage do we need?
\item What \textbf{level of performance} is required? 
\item What \textbf{method of access} is needed?
\item Is the storage \textbf{shared or exclusive} for the use of one VM?
\item How much does it \textbf{cost}?
\item How does it relate to \textbf{regions and availability zones}? 
\end{enumerate}

\section{Block storage}

Block storage is where the customer can provision emulated block devices, in other words virtual hard disk drives:
\begin{itemize}
\item As they emulate a hard disk, they can be attached to only one VM (at a time).
\item The VM is responsible for formatting, creating a filesystem (NTFS, ext3, etc) and managing access controls  on the emulated hard disk.
\item They can be resized \footnote{On AWS, not all providers allow this} but following an increase in size the filesystem must be expanded on the guest VM. 
\end{itemize}
When a compute instance is provisioned, its root volume is a block storage volume that is usually cloned from an initial image. This is often \textit{ephemeral} and disappears when the VM is destroyed.

\section{File storage}

File Storage is where the customer can provision shared file systems.
To the guest VM, this looks like a Network Drive.
\begin{itemize}
\item Multiple VMs can mount (i.e. access) the storage at one time.
\item Storage is accessed using file-sharing protocols as would be seen on internal LANs, such as NFS or SMB.
\item Primary access control is implemented by the provider according to customer settings.
\end{itemize}

\section{Object storage}

Object storage is where the customer can place objects (i.e. files) into managed storage using a key and access via REST-like web service.
\begin{itemize}
\item Storage can be accessed by VMs and over the internet using standard HTTP methods: \texttt{GET, PUT, DELETE} etc.
\item Needs no support on application other than ability to connect to HTTP-based web services.
\item Primary access control is implemented by the provider according to customer settings.
\item Application doesn't need to access its own guest VM's filesystem.
\item Doesn't usually have folders/directories (but can pretend to!)
\item Can serve content directly to internet as a public static website.
\end{itemize}

\section{Key question}

What is the correct type of storage (block, file or object) for a given application?

\end{document}
